% Pursuing Stacks (cotes 134-2 to 134-8): archive cote 134-7, site page 22 (PDF page 23), typescript page 464.
% typescript. Draft compared: ps6.tex:1023-1079.
% AI-transcribed by Claude Opus 5.5 (claude-opus-5-5, adaptive thinking, effort high), Message Batches API, batch ps3 (msgbatch_01XdFmuVVBAcgc7EFxdtCqsj), 2026-11-12; not yet checked by eye.
\scanpage{134-7/22}
\mspage{464}

\deleted{and these} should be submersions for a ``\uline{Kan complex}''. In case $X_*$ is
defined by a monoid object as above, this clearly implies that G is a
group \deleted{(surjectivity \unsure{of} \unsure{(xx)|(**)} on points should be enough, a fortiori it
would suffice that $X_*(k)$ is a Kan complex in the usual sense)} - which
rules out the counterexample above !

Of course, the very first thing we'll expect from a ``good notion'' of
Kan complexes in U(k), is that for $k=\underline{Z}$ , it should make the linearizati-%
on theorem work, namely the maps in (w) p.454 are quasi-isomorphisms.
The next thing, very close to this one but for arbitrary ground ring k
now, is that a map
\[ X_* \longrightarrow X'_* \qquad (\text{ with } X_0=X'_0=X_1=X'_1 = e \quad ) \]
of Kan complexes in U(k) is a homotopism, iff it induces an isomorphism
on the homology modules (z') (p.458) - which sounds reasonable precise
-ly because we are working with Kan complexes. If this is so, the
homotopy category $\mathrm{Hot}_1(k)$ of \xout{simply} 1-connected homotopy types over k
may be identified with \over{the}{a} category of Kan complexes \xout{as usual} ``up to
homotopy'' , as usual (but working \inserted{now} with complexes of unipotent bundles
over k). Third thing, still over arbitrary ground ring, would be
a development of the usual homotopy formalism in the unipotent context,
including (one hopes) homotopy fibers of maps, and P\over{\unsure{a}}{o}stnikov \xout{fibrations}
d\hand{\'e}vissage. Again, it is hard to imagine how to get such d\hand{\'e}vissage,
without getting hold inductively of homotopy invariants $\hand{$\pi$}_i$ which are
k-modules. This should come out if we are able to define homotopy fibers,
as for an (n-1)-connected k-homotopy type (defined here as one whose
homology invariants $H_i$ are zero for $i\hand{$\leqslant$}n-1$), $\hand{$\pi$}_n$ should be no more, no
less than $H_n$ , which is indeed a k-module. Coming back to $k=\underline{Z}$ again,
this should imply that \xout{the} for $n\hand{$\geqslant$}1$ at any rate, the canonical functor
\[ \hand{(***)} \qquad \mathrm{Hot}_n(\underline{Z}) \hand{$\longrightarrow$} \mathrm{Hot}_n \]
induces a bijection between isomorphism classes of schematic and ordinary
\inserted{n-connected} homotopy types - and it will be \over{\unsure{g}}{h}ard to believe this can be so, without
this functor being actually an equivalence of categories - the expected
apotheosis of the theory ! Maybe to this end, one may even be able to
introduce reasonable internal \uline{Hom}'s within the category of schematic
Kan complexes, \xout{\unsure{and show that they are}} \inserted[typed]{in a way} compatible with the familiar notion
in the discrete set-up.

\margin{22}
